% 12_mini_soc_de10lite.tex -- a RISC-V CPU + mm_core on one Avalon-MM bus in the MAX 10.
% Build: tectonic 12_mini_soc_de10lite.tex
\documentclass[tikz,border=10pt,12pt]{standalone}
% soc_style.tex -- shared look for the SoC figures (same palette as the slide decks and the
% core register placement figure). Every figure does % soc_style.tex -- shared look for the SoC figures (same palette as the slide decks and the
% core register placement figure). Every figure does % soc_style.tex -- shared look for the SoC figures (same palette as the slide decks and the
% core register placement figure). Every figure does \input{soc_style}.
\usepackage[T1]{fontenc}
\usepackage[scaled=0.95]{helvet}
\renewcommand{\familydefault}{\sfdefault}
\usepackage{booktabs}
\usepackage{amsmath}
\usepackage{amssymb}
\usetikzlibrary{arrows.meta,calc,positioning,fit,shapes.geometric,shapes.misc}

% ---------------- palette ----------------
\definecolor{pcb}{HTML}{0B2E33}
\definecolor{pcbmid}{HTML}{1F5C61}
\definecolor{tint}{HTML}{E6F0EF}
\definecolor{copper}{HTML}{B5651D}
\definecolor{coppert}{HTML}{F6E8DC}
\definecolor{ink}{HTML}{1F2A2E}
\definecolor{muted}{HTML}{5B6B70}
\definecolor{hardf}{HTML}{DCE6F2}
\definecolor{harde}{HTML}{2F4A7A}
\definecolor{good}{HTML}{2E7D32}
\definecolor{bad}{HTML}{B3261E}
% one colour per kind of block: border (d) / fill (f)
\definecolor{cpud}{HTML}{2B5C9E} \definecolor{cpuf}{HTML}{D6E4F5}
\definecolor{gpud}{HTML}{2E7D32} \definecolor{gpuf}{HTML}{D5EDD0}
\definecolor{npud}{HTML}{0B2E33} \definecolor{npuf}{HTML}{BFE0DB}
\definecolor{memd}{HTML}{5E3A8C} \definecolor{memf}{HTML}{E4DAF3}
\definecolor{busd}{HTML}{B5651D} \definecolor{busf}{HTML}{F6E8DC}
\definecolor{prd}{HTML}{546E7A}  \definecolor{prf}{HTML}{E3EAED}
\definecolor{secd}{HTML}{B3261E} \definecolor{secf}{HTML}{F7D9D6}
\definecolor{pwrd}{HTML}{8A6800} \definecolor{pwrf}{HTML}{FFF0B8}
\definecolor{medd}{HTML}{8E2463} \definecolor{medf}{HTML}{F5D6E6}
\definecolor{skd}{HTML}{2E7D32}  \definecolor{skf}{HTML}{BFE3B4}   % skew buffers, as in the register figure

\tikzset{
  blk/.style 2 args={draw=#1, fill=#2, rounded corners=4pt, thick, align=center,
                     font=\scriptsize, inner sep=3pt},
  cpu/.style={blk={cpud}{cpuf}},  gpu/.style={blk={gpud}{gpuf}},
  npu/.style={blk={npud}{npuf}},  mem/.style={blk={memd}{memf}},
  bus/.style={blk={busd}{busf}},  per/.style={blk={prd}{prf}},
  sec/.style={blk={secd}{secf}},  pwr/.style={blk={pwrd}{pwrf}},
  med/.style={blk={medd}{medf}},
  comb/.style ={draw=copper, fill=coppert, rounded corners=2pt, thick, align=center,
                font=\scriptsize, inner sep=2pt},                       % combinational logic
  hard/.style ={draw=harde, fill=hardf, thick, dashed, align=center,
                font=\scriptsize, inner sep=2pt},                       % hard macro / off-chip
  sram/.style ={draw=memd, fill=memf, thick, align=center, font=\scriptsize,
                inner sep=2pt},                                         % SRAM array, not flip-flops
  ff/.style   ={draw=#1, fill=#1!30, line width=0.35pt},                % one flip-flop square
  ffgone/.style={draw=muted, fill=white, densely dashed, line width=0.35pt},
  panel/.style={draw=muted, rounded corners=6pt, line width=0.9pt},
  ptitle/.style={anchor=north west, font=\normalsize\bfseries, text=ink, inner sep=0pt},
  note/.style ={font=\scriptsize, text=ink, align=left, inner sep=0pt},
  lab/.style  ={font=\tiny, text=muted, inner sep=1pt},
  axi/.style  ={{Stealth[length=2.6mm]}-{Stealth[length=2.6mm]}, line width=2.2pt, copper},
  axi1/.style ={-{Stealth[length=2.6mm]}, line width=2.2pt, copper},
  apb/.style  ={{Stealth[length=1.8mm]}-{Stealth[length=1.8mm]}, line width=1pt, copper},
  apb1/.style ={-{Stealth[length=1.8mm]}, line width=1pt, copper},
  dat/.style  ={-{Stealth[length=1.8mm]}, thick, pcbmid},
  ctl/.style  ={-{Stealth[length=1.8mm]}, densely dashed, thick, muted},
  irq/.style  ={-{Stealth[length=2mm]}, densely dotted, line width=1.3pt, bad},
  num/.style  ={circle, fill=copper, text=white, font=\scriptsize\bfseries, inner sep=1pt,
                minimum size=0.46cm},
  ours/.style ={draw=bad, line width=1.2pt, densely dashed, rounded corners=4pt},
}

% ---------------- helpers ----------------
\newcommand{\rn}[1]{\texttt{#1}}
% \ffrow{x}{y}{n}{cell}{style}: n flip-flops in a row, bottom-left corner at (x,y)
\newcommand{\ffrow}[5]{\foreach \q in {1,...,#3}{\draw[#5] ({#1+(\q-1)*#4},{#2}) rectangle ++(#4,#4);}}
% \ffcol{x}{y}{n}{cell}{style}: n flip-flops in a column, top-left corner at (x,y)
\newcommand{\ffcol}[5]{\foreach \q in {1,...,#3}{\draw[#5] ({#1},{#2-\q*#4}) rectangle ++(#4,#4);}}
% \ffgrid{x}{ytop}{n}{per row}{cell}{style}: n flip-flops, filled row by row from the top-left
\newcommand{\ffgrid}[6]{\pgfmathtruncatemacro{\gl}{#3-1}%
  \foreach \b in {0,...,\gl}{\pgfmathtruncatemacro{\gc}{mod(\b,#4)}\pgfmathtruncatemacro{\gr}{div(\b,#4)}%
    \draw[#6] ({#1+\gc*#5},{#2-(\gr+1)*#5}) rectangle ++(#5,#5);}}

% every figure: same page size, title, footer
\newcommand{\canvas}{\useasboundingbox (-0.3,-0.2) rectangle (38.6,30.4);}
\newcommand{\figtitle}[3]{%
  \node[num, minimum size=0.9cm, font=\Large\bfseries] at (0.45,29.55) {#1};
  \node[anchor=west, font=\LARGE\bfseries, text=pcb] at (1.15,29.55) {#2};
  \node[anchor=west, font=\small, text=ink] at (0,28.7) {#3};}
\newcommand{\figfoot}[1]{\node[anchor=west, font=\scriptsize, text=muted] at (0,0.25) {#1};}
% \panel{x0}{y0}{x1}{y1}{title}
\newcommand{\panel}[5]{\draw[panel] (#1,#2) rectangle (#3,#4);
  \node[ptitle] at ({#1+0.3},{#4-0.25}) {#5};}

% ---------------- timing diagrams: 1 cycle = \cw cm, low at y, high at y+0.45 ----------------
\newcommand{\cw}{1.0}
\newcommand{\wclk}[3]{% x0 y ncycles
  \foreach \k in {1,...,#3}{\draw[pcb, thick] ({#1+(\k-1)*\cw},#2) -- ++(0,0.45) -- ++(0.5*\cw,0)
     -- ++(0,-0.45) -- ++(0.5*\cw,0);}}
\newcommand{\wsig}[3]{% x0 y {level per cycle, e.g. 0,1,1,0}
  \gdef\wprev{x}%
  \foreach \v [count=\k from 0] in {#3}{%
    \draw[pcb, thick] ({#1+\k*\cw},{#2+0.45*\v}) -- ++(\cw,0);
    \if x\wprev\else\ifnum\wprev=\v\else\draw[pcb, thick] ({#1+\k*\cw},#2) -- ++(0,0.45);\fi\fi
    \global\let\wprev\v}}
\newcommand{\wbus}[5]{% x0 y c0 c1 text: one value held from cycle c0 to cycle c1
  \draw[pcbmid, thick, fill=white] ({#1+(#3)*\cw},{#2+0.225}) -- ({#1+(#3)*\cw+0.08},{#2+0.45})
    -- ({#1+(#4)*\cw-0.08},{#2+0.45}) -- ({#1+(#4)*\cw},{#2+0.225}) -- ({#1+(#4)*\cw-0.08},#2)
    -- ({#1+(#3)*\cw+0.08},#2) -- cycle;
  \node[font=\tiny\ttfamily] at ({#1+((#3)+(#4))*\cw/2},{#2+0.225}) {#5};}
\newcommand{\widle}[4]{% x0 y c0 c1: bus idle (no value)
  \draw[muted, thick] ({#1+(#3)*\cw},{#2+0.225}) -- ({#1+(#4)*\cw},{#2+0.225});}
\newcommand{\wname}[3]{\node[anchor=east, font=\scriptsize\ttfamily] at ({#1-0.15},{#2+0.225}) {#3};}
.
\usepackage[T1]{fontenc}
\usepackage[scaled=0.95]{helvet}
\renewcommand{\familydefault}{\sfdefault}
\usepackage{booktabs}
\usepackage{amsmath}
\usepackage{amssymb}
\usetikzlibrary{arrows.meta,calc,positioning,fit,shapes.geometric,shapes.misc}

% ---------------- palette ----------------
\definecolor{pcb}{HTML}{0B2E33}
\definecolor{pcbmid}{HTML}{1F5C61}
\definecolor{tint}{HTML}{E6F0EF}
\definecolor{copper}{HTML}{B5651D}
\definecolor{coppert}{HTML}{F6E8DC}
\definecolor{ink}{HTML}{1F2A2E}
\definecolor{muted}{HTML}{5B6B70}
\definecolor{hardf}{HTML}{DCE6F2}
\definecolor{harde}{HTML}{2F4A7A}
\definecolor{good}{HTML}{2E7D32}
\definecolor{bad}{HTML}{B3261E}
% one colour per kind of block: border (d) / fill (f)
\definecolor{cpud}{HTML}{2B5C9E} \definecolor{cpuf}{HTML}{D6E4F5}
\definecolor{gpud}{HTML}{2E7D32} \definecolor{gpuf}{HTML}{D5EDD0}
\definecolor{npud}{HTML}{0B2E33} \definecolor{npuf}{HTML}{BFE0DB}
\definecolor{memd}{HTML}{5E3A8C} \definecolor{memf}{HTML}{E4DAF3}
\definecolor{busd}{HTML}{B5651D} \definecolor{busf}{HTML}{F6E8DC}
\definecolor{prd}{HTML}{546E7A}  \definecolor{prf}{HTML}{E3EAED}
\definecolor{secd}{HTML}{B3261E} \definecolor{secf}{HTML}{F7D9D6}
\definecolor{pwrd}{HTML}{8A6800} \definecolor{pwrf}{HTML}{FFF0B8}
\definecolor{medd}{HTML}{8E2463} \definecolor{medf}{HTML}{F5D6E6}
\definecolor{skd}{HTML}{2E7D32}  \definecolor{skf}{HTML}{BFE3B4}   % skew buffers, as in the register figure

\tikzset{
  blk/.style 2 args={draw=#1, fill=#2, rounded corners=4pt, thick, align=center,
                     font=\scriptsize, inner sep=3pt},
  cpu/.style={blk={cpud}{cpuf}},  gpu/.style={blk={gpud}{gpuf}},
  npu/.style={blk={npud}{npuf}},  mem/.style={blk={memd}{memf}},
  bus/.style={blk={busd}{busf}},  per/.style={blk={prd}{prf}},
  sec/.style={blk={secd}{secf}},  pwr/.style={blk={pwrd}{pwrf}},
  med/.style={blk={medd}{medf}},
  comb/.style ={draw=copper, fill=coppert, rounded corners=2pt, thick, align=center,
                font=\scriptsize, inner sep=2pt},                       % combinational logic
  hard/.style ={draw=harde, fill=hardf, thick, dashed, align=center,
                font=\scriptsize, inner sep=2pt},                       % hard macro / off-chip
  sram/.style ={draw=memd, fill=memf, thick, align=center, font=\scriptsize,
                inner sep=2pt},                                         % SRAM array, not flip-flops
  ff/.style   ={draw=#1, fill=#1!30, line width=0.35pt},                % one flip-flop square
  ffgone/.style={draw=muted, fill=white, densely dashed, line width=0.35pt},
  panel/.style={draw=muted, rounded corners=6pt, line width=0.9pt},
  ptitle/.style={anchor=north west, font=\normalsize\bfseries, text=ink, inner sep=0pt},
  note/.style ={font=\scriptsize, text=ink, align=left, inner sep=0pt},
  lab/.style  ={font=\tiny, text=muted, inner sep=1pt},
  axi/.style  ={{Stealth[length=2.6mm]}-{Stealth[length=2.6mm]}, line width=2.2pt, copper},
  axi1/.style ={-{Stealth[length=2.6mm]}, line width=2.2pt, copper},
  apb/.style  ={{Stealth[length=1.8mm]}-{Stealth[length=1.8mm]}, line width=1pt, copper},
  apb1/.style ={-{Stealth[length=1.8mm]}, line width=1pt, copper},
  dat/.style  ={-{Stealth[length=1.8mm]}, thick, pcbmid},
  ctl/.style  ={-{Stealth[length=1.8mm]}, densely dashed, thick, muted},
  irq/.style  ={-{Stealth[length=2mm]}, densely dotted, line width=1.3pt, bad},
  num/.style  ={circle, fill=copper, text=white, font=\scriptsize\bfseries, inner sep=1pt,
                minimum size=0.46cm},
  ours/.style ={draw=bad, line width=1.2pt, densely dashed, rounded corners=4pt},
}

% ---------------- helpers ----------------
\newcommand{\rn}[1]{\texttt{#1}}
% \ffrow{x}{y}{n}{cell}{style}: n flip-flops in a row, bottom-left corner at (x,y)
\newcommand{\ffrow}[5]{\foreach \q in {1,...,#3}{\draw[#5] ({#1+(\q-1)*#4},{#2}) rectangle ++(#4,#4);}}
% \ffcol{x}{y}{n}{cell}{style}: n flip-flops in a column, top-left corner at (x,y)
\newcommand{\ffcol}[5]{\foreach \q in {1,...,#3}{\draw[#5] ({#1},{#2-\q*#4}) rectangle ++(#4,#4);}}
% \ffgrid{x}{ytop}{n}{per row}{cell}{style}: n flip-flops, filled row by row from the top-left
\newcommand{\ffgrid}[6]{\pgfmathtruncatemacro{\gl}{#3-1}%
  \foreach \b in {0,...,\gl}{\pgfmathtruncatemacro{\gc}{mod(\b,#4)}\pgfmathtruncatemacro{\gr}{div(\b,#4)}%
    \draw[#6] ({#1+\gc*#5},{#2-(\gr+1)*#5}) rectangle ++(#5,#5);}}

% every figure: same page size, title, footer
\newcommand{\canvas}{\useasboundingbox (-0.3,-0.2) rectangle (38.6,30.4);}
\newcommand{\figtitle}[3]{%
  \node[num, minimum size=0.9cm, font=\Large\bfseries] at (0.45,29.55) {#1};
  \node[anchor=west, font=\LARGE\bfseries, text=pcb] at (1.15,29.55) {#2};
  \node[anchor=west, font=\small, text=ink] at (0,28.7) {#3};}
\newcommand{\figfoot}[1]{\node[anchor=west, font=\scriptsize, text=muted] at (0,0.25) {#1};}
% \panel{x0}{y0}{x1}{y1}{title}
\newcommand{\panel}[5]{\draw[panel] (#1,#2) rectangle (#3,#4);
  \node[ptitle] at ({#1+0.3},{#4-0.25}) {#5};}

% ---------------- timing diagrams: 1 cycle = \cw cm, low at y, high at y+0.45 ----------------
\newcommand{\cw}{1.0}
\newcommand{\wclk}[3]{% x0 y ncycles
  \foreach \k in {1,...,#3}{\draw[pcb, thick] ({#1+(\k-1)*\cw},#2) -- ++(0,0.45) -- ++(0.5*\cw,0)
     -- ++(0,-0.45) -- ++(0.5*\cw,0);}}
\newcommand{\wsig}[3]{% x0 y {level per cycle, e.g. 0,1,1,0}
  \gdef\wprev{x}%
  \foreach \v [count=\k from 0] in {#3}{%
    \draw[pcb, thick] ({#1+\k*\cw},{#2+0.45*\v}) -- ++(\cw,0);
    \if x\wprev\else\ifnum\wprev=\v\else\draw[pcb, thick] ({#1+\k*\cw},#2) -- ++(0,0.45);\fi\fi
    \global\let\wprev\v}}
\newcommand{\wbus}[5]{% x0 y c0 c1 text: one value held from cycle c0 to cycle c1
  \draw[pcbmid, thick, fill=white] ({#1+(#3)*\cw},{#2+0.225}) -- ({#1+(#3)*\cw+0.08},{#2+0.45})
    -- ({#1+(#4)*\cw-0.08},{#2+0.45}) -- ({#1+(#4)*\cw},{#2+0.225}) -- ({#1+(#4)*\cw-0.08},#2)
    -- ({#1+(#3)*\cw+0.08},#2) -- cycle;
  \node[font=\tiny\ttfamily] at ({#1+((#3)+(#4))*\cw/2},{#2+0.225}) {#5};}
\newcommand{\widle}[4]{% x0 y c0 c1: bus idle (no value)
  \draw[muted, thick] ({#1+(#3)*\cw},{#2+0.225}) -- ({#1+(#4)*\cw},{#2+0.225});}
\newcommand{\wname}[3]{\node[anchor=east, font=\scriptsize\ttfamily] at ({#1-0.15},{#2+0.225}) {#3};}
.
\usepackage[T1]{fontenc}
\usepackage[scaled=0.95]{helvet}
\renewcommand{\familydefault}{\sfdefault}
\usepackage{booktabs}
\usepackage{amsmath}
\usepackage{amssymb}
\usetikzlibrary{arrows.meta,calc,positioning,fit,shapes.geometric,shapes.misc}

% ---------------- palette ----------------
\definecolor{pcb}{HTML}{0B2E33}
\definecolor{pcbmid}{HTML}{1F5C61}
\definecolor{tint}{HTML}{E6F0EF}
\definecolor{copper}{HTML}{B5651D}
\definecolor{coppert}{HTML}{F6E8DC}
\definecolor{ink}{HTML}{1F2A2E}
\definecolor{muted}{HTML}{5B6B70}
\definecolor{hardf}{HTML}{DCE6F2}
\definecolor{harde}{HTML}{2F4A7A}
\definecolor{good}{HTML}{2E7D32}
\definecolor{bad}{HTML}{B3261E}
% one colour per kind of block: border (d) / fill (f)
\definecolor{cpud}{HTML}{2B5C9E} \definecolor{cpuf}{HTML}{D6E4F5}
\definecolor{gpud}{HTML}{2E7D32} \definecolor{gpuf}{HTML}{D5EDD0}
\definecolor{npud}{HTML}{0B2E33} \definecolor{npuf}{HTML}{BFE0DB}
\definecolor{memd}{HTML}{5E3A8C} \definecolor{memf}{HTML}{E4DAF3}
\definecolor{busd}{HTML}{B5651D} \definecolor{busf}{HTML}{F6E8DC}
\definecolor{prd}{HTML}{546E7A}  \definecolor{prf}{HTML}{E3EAED}
\definecolor{secd}{HTML}{B3261E} \definecolor{secf}{HTML}{F7D9D6}
\definecolor{pwrd}{HTML}{8A6800} \definecolor{pwrf}{HTML}{FFF0B8}
\definecolor{medd}{HTML}{8E2463} \definecolor{medf}{HTML}{F5D6E6}
\definecolor{skd}{HTML}{2E7D32}  \definecolor{skf}{HTML}{BFE3B4}   % skew buffers, as in the register figure

\tikzset{
  blk/.style 2 args={draw=#1, fill=#2, rounded corners=4pt, thick, align=center,
                     font=\scriptsize, inner sep=3pt},
  cpu/.style={blk={cpud}{cpuf}},  gpu/.style={blk={gpud}{gpuf}},
  npu/.style={blk={npud}{npuf}},  mem/.style={blk={memd}{memf}},
  bus/.style={blk={busd}{busf}},  per/.style={blk={prd}{prf}},
  sec/.style={blk={secd}{secf}},  pwr/.style={blk={pwrd}{pwrf}},
  med/.style={blk={medd}{medf}},
  comb/.style ={draw=copper, fill=coppert, rounded corners=2pt, thick, align=center,
                font=\scriptsize, inner sep=2pt},                       % combinational logic
  hard/.style ={draw=harde, fill=hardf, thick, dashed, align=center,
                font=\scriptsize, inner sep=2pt},                       % hard macro / off-chip
  sram/.style ={draw=memd, fill=memf, thick, align=center, font=\scriptsize,
                inner sep=2pt},                                         % SRAM array, not flip-flops
  ff/.style   ={draw=#1, fill=#1!30, line width=0.35pt},                % one flip-flop square
  ffgone/.style={draw=muted, fill=white, densely dashed, line width=0.35pt},
  panel/.style={draw=muted, rounded corners=6pt, line width=0.9pt},
  ptitle/.style={anchor=north west, font=\normalsize\bfseries, text=ink, inner sep=0pt},
  note/.style ={font=\scriptsize, text=ink, align=left, inner sep=0pt},
  lab/.style  ={font=\tiny, text=muted, inner sep=1pt},
  axi/.style  ={{Stealth[length=2.6mm]}-{Stealth[length=2.6mm]}, line width=2.2pt, copper},
  axi1/.style ={-{Stealth[length=2.6mm]}, line width=2.2pt, copper},
  apb/.style  ={{Stealth[length=1.8mm]}-{Stealth[length=1.8mm]}, line width=1pt, copper},
  apb1/.style ={-{Stealth[length=1.8mm]}, line width=1pt, copper},
  dat/.style  ={-{Stealth[length=1.8mm]}, thick, pcbmid},
  ctl/.style  ={-{Stealth[length=1.8mm]}, densely dashed, thick, muted},
  irq/.style  ={-{Stealth[length=2mm]}, densely dotted, line width=1.3pt, bad},
  num/.style  ={circle, fill=copper, text=white, font=\scriptsize\bfseries, inner sep=1pt,
                minimum size=0.46cm},
  ours/.style ={draw=bad, line width=1.2pt, densely dashed, rounded corners=4pt},
}

% ---------------- helpers ----------------
\newcommand{\rn}[1]{\texttt{#1}}
% \ffrow{x}{y}{n}{cell}{style}: n flip-flops in a row, bottom-left corner at (x,y)
\newcommand{\ffrow}[5]{\foreach \q in {1,...,#3}{\draw[#5] ({#1+(\q-1)*#4},{#2}) rectangle ++(#4,#4);}}
% \ffcol{x}{y}{n}{cell}{style}: n flip-flops in a column, top-left corner at (x,y)
\newcommand{\ffcol}[5]{\foreach \q in {1,...,#3}{\draw[#5] ({#1},{#2-\q*#4}) rectangle ++(#4,#4);}}
% \ffgrid{x}{ytop}{n}{per row}{cell}{style}: n flip-flops, filled row by row from the top-left
\newcommand{\ffgrid}[6]{\pgfmathtruncatemacro{\gl}{#3-1}%
  \foreach \b in {0,...,\gl}{\pgfmathtruncatemacro{\gc}{mod(\b,#4)}\pgfmathtruncatemacro{\gr}{div(\b,#4)}%
    \draw[#6] ({#1+\gc*#5},{#2-(\gr+1)*#5}) rectangle ++(#5,#5);}}

% every figure: same page size, title, footer
\newcommand{\canvas}{\useasboundingbox (-0.3,-0.2) rectangle (38.6,30.4);}
\newcommand{\figtitle}[3]{%
  \node[num, minimum size=0.9cm, font=\Large\bfseries] at (0.45,29.55) {#1};
  \node[anchor=west, font=\LARGE\bfseries, text=pcb] at (1.15,29.55) {#2};
  \node[anchor=west, font=\small, text=ink] at (0,28.7) {#3};}
\newcommand{\figfoot}[1]{\node[anchor=west, font=\scriptsize, text=muted] at (0,0.25) {#1};}
% \panel{x0}{y0}{x1}{y1}{title}
\newcommand{\panel}[5]{\draw[panel] (#1,#2) rectangle (#3,#4);
  \node[ptitle] at ({#1+0.3},{#4-0.25}) {#5};}

% ---------------- timing diagrams: 1 cycle = \cw cm, low at y, high at y+0.45 ----------------
\newcommand{\cw}{1.0}
\newcommand{\wclk}[3]{% x0 y ncycles
  \foreach \k in {1,...,#3}{\draw[pcb, thick] ({#1+(\k-1)*\cw},#2) -- ++(0,0.45) -- ++(0.5*\cw,0)
     -- ++(0,-0.45) -- ++(0.5*\cw,0);}}
\newcommand{\wsig}[3]{% x0 y {level per cycle, e.g. 0,1,1,0}
  \gdef\wprev{x}%
  \foreach \v [count=\k from 0] in {#3}{%
    \draw[pcb, thick] ({#1+\k*\cw},{#2+0.45*\v}) -- ++(\cw,0);
    \if x\wprev\else\ifnum\wprev=\v\else\draw[pcb, thick] ({#1+\k*\cw},#2) -- ++(0,0.45);\fi\fi
    \global\let\wprev\v}}
\newcommand{\wbus}[5]{% x0 y c0 c1 text: one value held from cycle c0 to cycle c1
  \draw[pcbmid, thick, fill=white] ({#1+(#3)*\cw},{#2+0.225}) -- ({#1+(#3)*\cw+0.08},{#2+0.45})
    -- ({#1+(#4)*\cw-0.08},{#2+0.45}) -- ({#1+(#4)*\cw},{#2+0.225}) -- ({#1+(#4)*\cw-0.08},#2)
    -- ({#1+(#3)*\cw+0.08},#2) -- cycle;
  \node[font=\tiny\ttfamily] at ({#1+((#3)+(#4))*\cw/2},{#2+0.225}) {#5};}
\newcommand{\widle}[4]{% x0 y c0 c1: bus idle (no value)
  \draw[muted, thick] ({#1+(#3)*\cw},{#2+0.225}) -- ({#1+(#4)*\cw},{#2+0.225});}
\newcommand{\wname}[3]{\node[anchor=east, font=\scriptsize\ttfamily] at ({#1-0.15},{#2+0.225}) {#3};}

\begin{document}
\begin{tikzpicture}[line join=round]
\canvas
\figtitle{12}{Your next step: a mini SoC on the DE10-Lite (RISC-V CPU + your accelerator)}
  {Put a CPU in the same FPGA as \rn{mm\_core} and connect both to one bus. The CPU writes \rn{mm\_regs} directly:
   no PC and no UART in the data path. This is a plan, not something built yet.}

% =====================================================================
% (A) system diagram
% =====================================================================
\panel{0}{13.0}{24.0}{28.1}{The system in Platform Designer (Intel's SoC builder)}
\draw[pcb, line width=1.4pt, rounded corners=6pt] (0.35,15.2) rectangle (23.65,27.3);
\node[anchor=north west, font=\scriptsize\bfseries, text=pcb] at (0.5,27.2) {MAX 10 10M50 FPGA};
\node[pwr, minimum width=3.6cm, minimum height=1.3cm, text width=3.4cm, anchor=north west, font=\tiny] at (0.6,26.7)
  {\textbf{clock + reset}\\50 MHz board clock\\reset bridge (sync release)};
\node[cpu, minimum width=7.6cm, minimum height=2.1cm, text width=7.4cm, anchor=north west, font=\scriptsize] (cpu) at (4.6,26.7)
  {\textbf{RISC-V CPU}: Nios V/m (Intel's RV32I soft core) \textbf{or} your own RV32I core from the roadmap\\[2pt]
   \tiny instruction port + data port = 2 Avalon-MM hosts; IRQ inputs};
\node[npu, minimum width=10.6cm, minimum height=2.1cm, text width=10.4cm, anchor=north west, font=\scriptsize] at (12.6,26.7)
  {\textbf{why}: the CPU runs the control code (load A, B, START, wait, read C) at bus speed; the PC is only needed
   for printing results};
\draw[bus] (0.6,22.7) rectangle (23.4,23.9);
\node[font=\scriptsize\bfseries, text=busd] at (12.0,23.3)
  {\textbf{Avalon-MM interconnect}, generated: address decoders, arbiters, width adapters, waitrequest logic};
\draw[axi] (6.5,24.6) -- (6.5,23.9); \draw[axi] (10.3,24.6) -- (10.3,23.9);
\node[lab, text=ink, anchor=west] at (6.6,24.3) {instr}; \node[lab, text=ink, anchor=west] at (10.4,24.3) {data};
\foreach \nm/\tx/\sty [count=\k from 0] in {
    on-chip RAM/{M9K blocks, 32 KB: program + data}/mem,
    JTAG UART/{printf to the PC over the USB-Blaster}/per,
    interval timer/{counts cycles; OS tick}/per,
    PIO/{LEDs, switches, 7-segment}/per,
    SDRAM controller/{to the 64 MB SDRAM chip}/mem,
    \rn{mm\_core\_avalon}/{wrapper around your \rn{mm\_core}}/npu}{
  \pgfmathsetmacro{\x}{0.6+\k*3.82}
  \node[\sty, fill=white, minimum width=3.6cm, minimum height=2.3cm, text width=3.4cm, anchor=north west, font=\tiny]
    (ag\k) at (\x,21.9) {\textbf{\nm}\\[2pt]\tx};
  \draw[apb] ({\x+1.8},22.7) -- ({\x+1.8},21.9);}
\draw[ours] (19.65,22.0) rectangle (23.35,19.5);
\node[hard, minimum width=3.6cm, minimum height=1.1cm, text width=3.4cm, anchor=north west, font=\tiny] at (15.88,18.4)
  {\textbf{SDRAM chip}\\on the board, 16 bit};
\draw[axi] (17.68,19.6) -- (17.68,18.4);
\node[npu, minimum width=3.6cm, minimum height=1.6cm, text width=3.4cm, anchor=north west, font=\tiny] at (19.7,18.4)
  {\textbf{mm\_core}\\(unchanged: regs, RAM banks, skew, 4$\times$4 PEs, FSM)};
\draw[dat] (21.5,19.6) -- (21.5,18.4);
\draw[irq] (21.5,22.0) -- (21.5,22.4) -- (12.1,22.4) -- (12.1,24.6);
\draw[irq] (9.94,21.9) -- (9.94,22.4);
\node[font=\tiny\bfseries, text=bad, fill=busf, inner sep=1pt] at (15.0,22.4) {IRQ: DONE};
\node[note, anchor=north west, text width=15.0cm] at (0.6,18.35) {%
  The UART bridge, \rn{uart\_rx} and \rn{uart\_tx} are no longer in the data path. The JTAG UART is only for messages.\\[3pt]
  Platform Designer writes all the bus logic; you write one component: the wrapper (right).};

% =====================================================================
% (B) memory map
% =====================================================================
\panel{24.3}{20.6}{38.3}{28.1}{Memory map (example)}
\node[anchor=north west, font=\scriptsize, inner sep=0pt] at (24.6,27.35) {%
  \renewcommand{\arraystretch}{1.15}%
  \begin{tabular}{@{}l l l@{}}
    \toprule
    \textbf{Component} & \textbf{Base address} & \textbf{Size}\\
    \midrule
    SDRAM          & \rn{0x0000\_0000} & 64 MB\\
    on-chip RAM    & \rn{0x0800\_0000} & 32 KB\\
    \rn{mm\_core}  & \rn{0x0810\_0000} & 16 KB (your 0x0000--0x3FFF map)\\
    JTAG UART      & \rn{0x0811\_0000} & 8 B\\
    timer          & \rn{0x0811\_0020} & 32 B\\
    PIO (LEDs)     & \rn{0x0811\_0040} & 16 B\\
    \bottomrule
  \end{tabular}};
\node[note, font=\scriptsize, anchor=north west, text width=13.4cm] at (24.6,23.45) {%
  Platform Designer assigns the real base addresses; the generated \rn{system.h} gives them to your C code
  (e.g. \rn{MM\_CORE\_BASE}). Your register offsets stay the same: START is at base + 0x0000.};

% =====================================================================
% (C) the wrapper
% =====================================================================
\panel{24.3}{13.0}{38.3}{20.3}{The wrapper: Avalon-MM agent $\rightarrow$ your \rn{mm\_regs} bus}
\node[anchor=north west, font=\scriptsize, inner sep=0pt] at (24.6,19.55) {%
  \renewcommand{\arraystretch}{1.12}%
  \begin{tabular}{@{}l c l@{}}
    \toprule
    \textbf{Avalon agent signal} & & \textbf{your signal}\\
    \midrule
    \rn{avs\_read | avs\_write}   & $\rightarrow$ & \rn{bus\_valid}\\
    \rn{avs\_write}               & $\rightarrow$ & \rn{bus\_write}\\
    \rn{avs\_address}, \rn{avs\_byteenable} & $\rightarrow$ & \rn{bus\_addr} (byte address)\\
    \rn{avs\_writedata} (lane)    & $\rightarrow$ & \rn{bus\_wdata}\\
    \rn{avs\_waitrequest}         & $\leftarrow$  & \rn{!bus\_ready}\\
    \rn{avs\_readdata}            & $\leftarrow$  & \rn{bus\_rdata}\\
    \rn{avs\_readdatavalid}       & $\leftarrow$  & \rn{bus\_rvalid} (read latency 1)\\
    \rn{irq}                      & $\leftarrow$  & DONE flag AND an enable bit (new)\\
    \bottomrule
  \end{tabular}};
\node[note, font=\scriptsize, anchor=north west, text width=13.4cm] at (24.6,15.2) {%
  \textbf{Watch the byte lanes}: on a 32-bit bus a byte store (\rn{sb}) to A at offset 1 arrives on \rn{writedata[15:8]} with
  \rn{byteenable = 0010}. Your \rn{mm\_regs} takes A / B bytes from \rn{wdata[7:0]}, so the wrapper must pick the lane that
  \rn{byteenable} marks (and rebuild the low 2 address bits from it if the address arrives word-aligned).};

% =====================================================================
% (D) C code
% =====================================================================
\panel{0}{0.7}{12.6}{12.7}{The software: plain C on the RISC-V}
\node[draw=cpud, fill=white, rounded corners=3pt, anchor=north west, align=left, font=\scriptsize\ttfamily, inner sep=5pt] at (0.35,11.95) {%
  \#include "system.h"\ \ \ \ \ \ \ \ \ \ // MM\_CORE\_BASE from Platform Designer\\
  volatile uint8\_t\ \ *mm8\ \ = (uint8\_t *)\ MM\_CORE\_BASE;\\
  volatile uint32\_t *mm32 = (uint32\_t *)MM\_CORE\_BASE;\\[3pt]
  for (i = 0; i < 4; i++)\ \ \ \ \ \ \ \ \ \ \ \ // A: 0x1000 + i*64 + k\\
  \ \ for (k = 0; k < K; k++) mm8[0x1000 + i*64 + k] = A[i][k];\\
  for (k = 0; k < K; k++)\ \ \ \ \ \ \ \ \ \ \ \ // B: 0x2000 + k*4 + j\\
  \ \ for (j = 0; j < 4; j++) mm8[0x2000 + k*4 + j] = B[k][j];\\[3pt]
  mm32[0x0008/4] = K;\ \ \ \ \ \ \ \ \ \ \ \ \ \ \ \ // CFG\_K\\
  mm32[0x0000/4] = 1;\ \ \ \ \ \ \ \ \ \ \ \ \ \ \ \ // CTRL.START\\
  while ((mm32[0x0004/4] \& 2) == 0) ;\ // STATUS.DONE (or wait for IRQ)\\[3pt]
  for (i = 0; i < 16; i++)\ \ \ \ \ \ \ \ \ \ \ // C: 0x3000 + 4*index\\
  \ \ C[i] = (int32\_t)mm32[(0x3000 + 4*i)/4];};
\node[note, anchor=north west, text width=11.9cm] at (0.35,6.6) {%
  Exactly the same register map that \rn{host.py} uses today, just written as memory stores and loads.
  \rn{volatile} tells the compiler that every access must really happen on the bus.\\[3pt]
  If your CPU has a data cache, keep this address range uncached, or flush (Figure 8). Nios V/m has no data cache.};

% =====================================================================
% (E) speed comparison
% =====================================================================
\panel{12.9}{0.7}{25.4}{12.7}{Where the time goes: today vs mini SoC}
\foreach \nm/\v/\col/\txt [count=\k from 0] in {
    {compute (12 cycles)}/0.38/npud/{240 ns (measured in simulation)},
    {mini SoC: bus writes + reads}/1.88/cpud/{$\approx$5--10 $\mu$s (estimate)},
    {today: 154 UART bytes}/5.13/bad/{13.4 ms}}{
  \pgfmathsetmacro{\yy}{10.6-\k*1.6}
  \node[font=\scriptsize\bfseries, anchor=west] at (13.2,{\yy+0.55}) {\nm};
  \fill[\col!55] (13.2,{\yy-0.25}) rectangle ++({0.3+\v*1.9},0.5);
  \node[font=\scriptsize, anchor=west] at ({13.65+\v*1.9},\yy) {\txt};}
\draw[muted] (13.2,5.85) -- (24.9,5.85);
\foreach \t/\x in {100 ns/0, 1 $\mu$s/1, 10 $\mu$s/2, 100 $\mu$s/3, 1 ms/4, 10 ms/5}{
  \draw[muted] ({13.5+\x*1.9},5.75) -- ({13.5+\x*1.9},5.95);
  \node[lab] at ({13.5+\x*1.9},5.55) {\t};}
\node[lab, text=ink, anchor=north] at (19.0,5.3) {log scale: each step is 10$\times$};
\node[note, anchor=north west, text width=11.9cm] at (13.2,4.5) {%
  Today the UART moves 154 bytes at 86.8 $\mu$s each: about 56,000$\times$ the compute time.
  On the bus, about 50 transfers of a few CPU cycles each: roughly 1,000$\times$ less overhead.\\[3pt]
  The estimate assumes a simple CPU at 50 MHz; measure the real number with the interval timer and put it in your README.};

% =====================================================================
% (F) steps
% =====================================================================
\panel{25.7}{0.7}{38.3}{12.7}{Build steps (on Ubuntu)}
\node[note, anchor=north west, text width=12.0cm] at (26.0,11.95) {%
  \textbf{1} Platform Designer: new system; add the CPU, on-chip RAM, JTAG UART, timer, PIO.
  Check that your Quartus Lite version lists Nios V/m for MAX 10 in the IP Catalog.\\[3pt]
  \textbf{2} Write \rn{mm\_core\_avalon.sv} (the wrapper) and add it as a component with the Component Editor.\\[3pt]
  \textbf{3} Add the IRQ output: DONE AND an enable bit.\\[3pt]
  \textbf{4} Generate the system, instantiate it in \rn{top\_de10lite.sv}, compile, check timing at 50 MHz.\\[3pt]
  \textbf{5} Make the board support package + app (Nios V command-line tools or the RiscFree IDE).\\[3pt]
  \textbf{6} Run; print C through the JTAG UART; compare with \rn{golden\_model.py}.\\[3pt]
  \textbf{7} Time it with the interval timer; record the numbers.\\[5pt]
  \textbf{Other CPU options}: your own RV32I core (roadmap step 2), or an open core such as PicoRV32 with a small
  bus adapter. The wrapper idea stays the same.};

\figfoot{Base addresses, sizes and the overhead estimate are examples / estimates. Tool and IP availability depends on the Quartus version: check the IP Catalog.}
\end{tikzpicture}
\end{document}
